%============================ INTRODUCCIÓN ==================================
\chapter{Introducción}
\label{cap:introduccion}

La banda de 2.4~GHz constituye hoy uno de los recursos radioeléctricos más congestionados del entorno cotidiano. En ella coexisten redes Wi\nobreakdash-Fi bajo el estándar IEEE~802.11, dispositivos Bluetooth y Bluetooth de Baja Energía (BLE), teléfonos inalámbricos, hornos de microondas y una densidad creciente de dispositivos del Internet de las Cosas (IoT). Esta coexistencia, sumada a que solo tres de los trece canales de 2.4~GHz (1, 6 y 11) no se solapan entre sí, provoca degradaciones de rendimiento difíciles de diagnosticar sin instrumentación especializada \parencite{gast2005,bianchi2000}. Los analizadores de espectro comerciales, si bien precisos, resultan costosos y poco accesibles para pequeños negocios, instituciones educativas y laboratorios de aprendizaje.

Paralelamente, la seguridad de las redes inalámbricas se ha vuelto una preocupación central. Ataques como el punto de acceso \emph{gemelo malicioso} (\emph{evil twin}), en el que un atacante replica el nombre de una red legítima para interceptar tráfico, explotan precisamente la naturaleza abierta del medio y la confianza del usuario \parencite{callegati2009,owasp2021}. Detectar la presencia de estos puntos de acceso y comprender la ocupación del espectro son capacidades formativas esenciales para quien se inicia en la ciberseguridad.

El proyecto \textbf{EILOR} nace en la intersección de ambas necesidades: ofrecer una plataforma \emph{asequible, didáctica y funcional} que permita, por un lado, \emph{visualizar y analizar} el espectro de 2.4~GHz en tiempo real y, por otro, ejercer un \emph{control de acceso} responsable mediante un punto de acceso propio con portal cautivo de registro de invitados ---al estilo del Wi\nobreakdash-Fi de una cafetería---. Todo ello se construye sobre hardware de bajo costo (ESP32) y una pila de software abierta (Node.js, Express, WebSocket).

\section{Justificación}
La elección del microcontrolador ESP32 responde a que integra, en un mismo encapsulado y por un costo reducido, radios Wi\nobreakdash-Fi y Bluetooth/BLE, capacidad de procesamiento suficiente para escaneo activo y una pila TCP/IP completa \parencite{espressif2023}. Frente a soluciones cerradas, EILOR privilegia la \emph{transparencia} (código propio y modular), la \emph{reproducibilidad} (persistencia en archivos JSON sin dependencias nativas) y el \emph{valor pedagógico}: cada componente ---desde el cálculo de interferencia por canal hasta la autenticación por token--- es observable y auditable.

\section{Alcance y delimitación}
El sistema se circunscribe a la banda de 2.4~GHz y a los canales de \emph{advertising} de BLE (37, 38 y 39). No abarca la banda de 5~GHz ni la decodificación de tramas de datos de terceros; el escaneo Wi\nobreakdash-Fi se limita a la información de baliza (SSID, BSSID, canal, RSSI y tipo de cifrado) que la propia API del ESP32 expone. La funcionalidad de punto de acceso y portal cautivo está concebida para operar sobre una \emph{red propia}, con un nombre de red identificable y un aviso explícito de registro, en cumplimiento de principios éticos y de protección de datos; su uso para suplantar redes ajenas queda expresamente fuera del alcance autorizado del proyecto.

\section{Estructura del documento}
El documento integra los cuatro entregables de la macro actividad. El Capítulo~\ref{cap:problema} delimita el problema en sus niveles macro, meso y micro; el Capítulo~\ref{cap:marco} desarrolla el marco teórico; el Capítulo~\ref{cap:objetivos} formula los objetivos bajo criterios SMART; y el Capítulo~\ref{cap:arquitectura} presenta la arquitectura y las tecnologías (Entrega~1). El Capítulo~\ref{cap:diseno} detalla el diseño; el Capítulo~\ref{cap:implementacion} documenta la implementación con el código real del sistema; y el Capítulo~\ref{cap:pruebas} expone el plan y los resultados de pruebas de seguridad (Entrega~2). Los Capítulos~\ref{cap:desempeno} a \ref{cap:manuales} abordan el desempeño, el análisis económico\nobreakdash-ambiental con curva~S, la escalabilidad y los manuales (Entrega~3). Finalmente, el Capítulo~\ref{cap:conclusiones} presenta las conclusiones y la sustentación (Entrega~4), seguido de las referencias y los anexos.
